\vspace{-4pt}
\section{Experiments}\label{sec:experiments}

% Evaluation Protocol Figure (floats here so it lands near its first reference in section 3.6)
\begin{figure}[!t]
  \centering
  \includegraphics[width=0.85\linewidth]{figures/assets/evaluation-protocol.pdf}
  \caption{Flow aware evaluation runtime.}
  \label{fig:protocol}
\end{figure}

% --- Macaron palette for RQ1 design tables ---
\definecolor{macaronPinkHead}{HTML}{F4B6C2}
\definecolor{macaronPinkBody}{HTML}{FCE4EC}
\definecolor{macaronBlueHead}{HTML}{B6D4F4}
\definecolor{macaronBlueBody}{HTML}{E3F0FC}
\definecolor{macaronMintHead}{HTML}{B7E4C7}
\definecolor{macaronMintBody}{HTML}{E6F5EC}
\definecolor{macaronLemonHead}{HTML}{F7E29A}
\definecolor{macaronLemonBody}{HTML}{FBF3D0}
\definecolor{macaronLavHead}{HTML}{D4C2F0}
\definecolor{macaronLavBody}{HTML}{EEE7F8}

% --- Answer box (matches ARXIV/context-bench summary style) ---
\definecolor{lhbAnswerBg}{RGB}{243,248,252}
\newenvironment{summary}{%
  \begin{tcolorbox}[width=\linewidth, colback=lhbAnswerBg, top=1pt, bottom=1pt, left=2pt, right=2pt, before skip=3pt, after skip=3pt]%
}{%
  \end{tcolorbox}%
}

We organize the experiments around three questions that treat coding agent evaluation as loop engineering evaluation.
All experiments use the full 112 task release of \lhb{}.
\vspace{-3pt}
\subsection{RQ1: How Far Do Coding Agent Loops Sustain Progress?}\label{sec:rq1}

\autoref{tab:rq1} reports RQ1 results along three dimensions that separate model choice, loop implementation, and residual continuation. The table covers within family model scaling under the corresponding vendor loop, model sweeps under a fixed loop implementation, and loop implementation sweeps under a fixed model.
Across all three dimensions, every cell remains well below full resolution of \lhb{}, suggesting that long horizon limitations are not explained by model choice or loop implementation in isolation.
The strongest model and loop configuration resolves 25.00\% of tasks. Without the outer evaluation loop, a single evaluated execution segment usually reaches a shallow dependency prefix.
These pairings align frontier models with the vendor infrastructure under which they are typically deployed, yet long horizon loop execution remains well below complete resolution.
The loop implementation sweep isolates how the surrounding execution loop changes progress under the same model.
Without external continuation, evaluated loop implementations cover a shallow prefix under their own state retention and work routing policies.
The recorded loop traces also capture partial progress patterns, including unresolved units with empty patches and skipped obligations that prioritize lower burden units.
External continuation reduces early stalls and increases partial progress, while residual routing and state retention remain central loop engineering factors.
In the model dimension, Resolve Rate plateaus even when the outer evaluation loop attempts every unit, because later units inherit earlier open obligations.
Even under tool use and execution oriented training, the evaluated model and loop configurations often decompose the task into local issues without consistently maintaining a persistent global plan. Appendix~\ref{app:passk} reports pass@K to check that within family monotonicity is not driven by a single trial. In Table~\ref{tab:rq1}, RR denotes Resolve Rate, TPR denotes Test Pass Rate, and Depth is normalized dependency depth. Mean billed tokens per task appear in Appendix~\ref{app:rq1-tokens}.

\begin{table*}[!t]
  \centering
  \caption{RQ1 performance across models and loop implementations. Per task billed tokens are reported in Appendix~\ref{app:rq1-tokens}.}
  \label{tab:rq1}
  \begin{minipage}[t]{0.49\textwidth}
    \vspace{0pt}
    \centering
    \footnotesize
    \setlength{\tabcolsep}{5pt}
    \renewcommand{\arraystretch}{1.16}
    \resizebox{\linewidth}{!}{%
    \begin{tabular}{@{}l cc cc c@{}}
      \toprule
      \textbf{Model} & \multicolumn{2}{c}{\textbf{RR}$\uparrow$} & \multicolumn{2}{c}{\textbf{TPR}$\uparrow$} & \textbf{Depth}$\uparrow$ \\
       & \textit{w/o} & \textit{w/} & \textit{w/o} & \textit{w/} & \\
      \midrule
      \rowcolor{macaronMintHead}\multicolumn{6}{l}{\emph{GPT (Codex)}} \\
      GPT-5.5       & 14.29\% & 21.43\% & 39.62\% & 50.84\% & 0.53 \\
      GPT-5.2       & 9.82\%  & 13.39\% & 33.91\% & 43.90\% & 0.36 \\
      GPT-5         & 5.36\%  & 8.04\%  & 28.45\% & 36.86\% & 0.22 \\
      GPT-4o        & 1.79\%  & 3.57\%  & 21.85\% & 28.30\% & 0.09 \\
      \rowcolor{macaronMintHead}\multicolumn{6}{l}{\emph{Claude (Claude Code)}} \\
      Opus-4.7      & 16.96\% & 25.00\% & 41.18\% & 53.05\% & 0.61 \\
      Opus-4.5      & 12.50\% & 17.86\% & 36.97\% & 48.21\% & 0.46 \\
      Sonnet-4.5    & 8.93\%  & 12.50\% & 32.85\% & 42.58\% & 0.34 \\
      Sonnet-4      & 4.46\%  & 7.14\%  & 27.43\% & 35.54\% & 0.18 \\
      \rowcolor{macaronMintHead}\multicolumn{6}{l}{\emph{Qwen (Qwen Code)}} \\
      Qwen3.6-Plus  & 6.25\%  & 9.82\%  & 37.97\% & 45.10\% & 0.31 \\
      Qwen3.5-Plus  & 5.36\%  & 8.04\%  & 31.75\% & 41.17\% & 0.19 \\
      Qwen3-Max     & 0.89\%  & 1.79\%  & 17.45\% & 23.31\% & 0.07 \\
      Qwen2.5-72B   & 0.00\%  & 0.00\%  & 10.61\% & 13.88\% & 0.03 \\
      \noalign{\vskip 6.4pt}
      \bottomrule
    \end{tabular}%
    }
  \end{minipage}\hfill
  \begin{minipage}[t]{0.49\textwidth}
    \vspace{0pt}
    \centering
    \footnotesize
    \setlength{\tabcolsep}{5pt}
    \renewcommand{\arraystretch}{1.16}
    \resizebox{\linewidth}{!}{%
    \begin{tabular}{@{}l cc cc c@{}}
      \toprule
      & \multicolumn{2}{c}{\textbf{RR}$\uparrow$} & \multicolumn{2}{c}{\textbf{TPR}$\uparrow$} & \textbf{Depth}$\uparrow$ \\
       & \textit{w/o} & \textit{w/} & \textit{w/o} & \textit{w/} & \\
      \midrule
      \rowcolor{macaronPinkHead}\multicolumn{6}{l}{\textbf{Model} \textit{(fixed loop: Claude Code)}} \\
      Opus-4.7      & 16.96\% & 25.00\% & 41.18\% & 53.05\% & 0.61 \\
      GPT-5.5       & 13.39\% & 20.54\% & 38.04\% & 49.30\% & 0.47 \\
      GLM-5.1       & 13.39\% & 18.75\% & 34.55\% & 45.49\% & 0.44 \\
      DeepSeek-V4P  & 11.61\% & 18.75\% & 36.07\% & 47.34\% & 0.46 \\
      Gemini-3.1-Pro & 8.93\%  & 14.29\% & 32.86\% & 43.59\% & 0.31 \\
      Qwen3.6-Plus  & 6.25\%  & 9.82\%  & 36.43\% & 47.50\% & 0.34 \\
      Kimi-2.6      & 3.57\%  & 6.25\%  & 21.27\% & 28.33\% & 0.14 \\
      Grok-4.1-FR   & 2.68\%  & 4.46\%  & 23.32\% & 30.97\% & 0.15 \\
      \rowcolor{macaronBlueHead}\multicolumn{6}{l}{\textbf{Loop} \textit{(fixed model: \texttt{gpt-5.4})}} \\
      Codex          & 13.39\% & 18.75\% & 37.84\% & 49.06\% & 0.49 \\
      Claude Code    & 12.50\% & 17.86\% & 36.97\% & 48.21\% & 0.42 \\
      GitHub Copilot & 10.71\% & 15.18\% & 34.21\% & 45.46\% & 0.38 \\
      OpenHands      & 6.25\%  & 9.82\%  & 29.21\% & 39.30\% & 0.20 \\
      SWE-agent      & 5.36\%  & 8.93\%  & 28.07\% & 38.11\% & 0.16 \\
      mini-swe-agent & 4.46\%  & 7.14\%  & 25.96\% & 35.20\% & 0.13 \\
      \bottomrule
    \end{tabular}%
    }
  \end{minipage}
\end{table*}

\begin{summary}
\textbf{RQ1:} Sustained long horizon loop progress remains limited at frontier scale. The limitation appears in loop level obligations such as global plan maintenance, residual continuity, and regression obligation retention rather than in per unit code generation alone.
\end{summary}

\vspace{-3pt}
\subsection{RQ2: How Do Coding Agent Loops Maintain Plan, Code, and Test State?}\label{sec:rq2}
% RQ2 grid: loop trace profile vs human development.
\begin{table}[t]
  \centering
  \caption{Planning, implementation, and testing loop trace metrics.}
  \label{tab:rq2-grid}
  \definecolor{rqGroupRule}{HTML}{ECECEC}
  \tiny
  \setlength{\tabcolsep}{3.6pt}
  \renewcommand{\arraystretch}{1.05}
  \resizebox{\linewidth}{!}{%
  \begin{tabular}{@{}l cccc cc ccc @{}}
    \toprule
    & \multicolumn{4}{c}{\textbf{Planning fidelity}}
    & \multicolumn{2}{c}{\textbf{Implementation}}
    & \multicolumn{3}{c}{\textbf{Testing}} \\
    \cmidrule(lr){2-5} \cmidrule(lr){6-7} \cmidrule(lr){8-10}
    \textbf{Loop}
      & \textbf{Edge F1} & \textbf{Layer~$\rho$} & \textbf{CPR} & \textbf{WR}
      & \textbf{PatchLen} & \textbf{Jacc}
      & \textbf{\#T} & \textbf{F2P} & \textbf{Reg}$\downarrow$ \\
    \midrule
    \rowcolor{rqGroupRule}\multicolumn{10}{l}{\emph{Closed source loop implementations}} \\
    Claude Code
      & 0.71 & 0.65 & 0.31 & 1.08
      & 1.58 & 0.62
      & 28 & 0.47 & 7.11\% \\
    Codex
      & 0.67 & 0.61 & 0.33 & 1.14
      & 1.71 & 0.64
      & 24 & 0.44 & 4.83\% \\
    GitHub Copilot
      & 0.58 & 0.52 & 0.41 & 0.92
      & 1.83 & 0.66
      & 22 & 0.41 & 6.91\% \\
    \rowcolor{rqGroupRule}\multicolumn{10}{l}{\emph{Open source loop implementations}} \\
    OpenHands
      & 0.39 & 0.33 & 0.85 & 0.39
      & 2.31 & 0.73
      & 16 & 0.34 & 2.46\% \\
    SWE-agent
      & 0.37 & 0.31 & 0.88 & 0.36
      & 2.36 & 0.74
      & 15 & 0.33 & 2.18\% \\
    mini-swe-agent
      & 0.27 & 0.22 & 0.97 & 0.24
      & 2.54 & 0.76
      & 11 & 0.28 & 0.24\% \\
    \bottomrule
  \end{tabular}%
  }
\end{table}


\rev{\autoref{tab:rq2-grid} uses recorded loop traces to compare evaluated loops with the source recovered prerequisite DAG along planning, implementation, and testing axes. The runtime permits any topological execution order consistent with the DAG, while the planning diagnostics score source evidenced relations and layers rather than reproduction of a historical total order.}
\textbf{(1) Planning: an internal dependency gap.}
\rev{Recorded plans recover partial dependency structure for long horizon routing, omitting many source evidenced pull request level prerequisites in the reference DAG.}
\rev{Edge F1 and Layer~$\rho$ remain far from perfect agreement with the source recovered prerequisite DAG for every evaluated loop, with lower values for open loop implementations than for closed ones.}
Critical Path Ratio (CPR) and Width Ratio (WR) reveal how missing prerequisites reshape the planned concurrency budget.
\rev{Closed source loop implementations organize execution through a tree shaped concurrent structure and stay close to the concurrency budget implied by the reference DAG, whereas open source loop implementations driven by linear control flow form a near chain.}
\rev{Width Ratio above unity for closed loop implementations is better interpreted as overparallelization than as planning fidelity beyond the source recovered reference DAG.}
A limited share of those branches reflects valid alternative decomposition, while most parallelize serial requirements absent from the recorded plan.
\textbf{(2) Implementation: patch surplus over a comparable token base.}
On units that the evaluated configuration eventually resolves, PatchLen exceeds the gold reference by a consistent margin while token overlap with gold remains moderate. Candidate patches are therefore often longer than the reference while preserving enough token overlap to remain comparable to it. Appendix~\ref{app:rq2-extended} provides a diff level account of where the extra lines come from.
\textbf{(3) Testing: a maintenance gap over the long horizon.}
Recorded loop traces contain sparse test authoring as long horizon work unfolds, leaving earlier obligations with limited agent authored regression protection against later edits.
Test count (\#T) is markedly lower than that of the native suite for every evaluated loop. Fail to pass yield (F2P) on authored tests remains acceptable after internal validation before commit. Regression (Reg) captures the downstream cost of this sparse maintenance.
Deployed loop implementations exhibit several percentage points of regression over previously satisfied obligations. The lower apparent Regression Rate of lighter implementations reflects the limited number of passing obligations eligible for regression measurement.
Long horizon loop engineering evaluation derives loop trace metrics from the running state of the evaluated loop, recording where plan, code, and test state are maintained, rerouted, or degraded.
\begin{summary}
\textbf{RQ2:} The long horizon gap centers on loop state discipline. Recorded loop traces capture missing dependencies, inflated patches, and sparse tests, linking performance to how plan, code, and test state are maintained across the run.
\end{summary}

\vspace{-3pt}
\subsection{RQ3: Which Loop Engineering Factors Shape Long Horizon Execution?}\label{sec:rq3}

RQ3 studies objective persistence, residual routing, context renewal, and regression pressure. We compare four representative loop runs. \emph{Codex goal mode} keeps a persistent objective inside the Codex loop~\citep{openai_codex_prompting}. \emph{Claude Code goal mode} maintains a session condition checked after each turn~\citep{anthropic_claude_code_goal}. \emph{Claude dynamic workflows} launch task specific workers with narrower local contexts~\citep{anthropic_claude_code_workflows,anthropic_claude_code_subagents}. The Ralph loop starts a fresh invocation for residual work~\citep{huntley_ralph_wiggum,huntley_ralph_loop}.
\autoref{tab:rq3-loop-engineering} reports context-budget rounds, regression events per run, and Resolve Rate. A context-budget round is the initial context or a segment opened by the implementation compaction policy or by an external restart. \autoref{fig:rq3-compaction-ablation} shows the compaction-budget sensitivity. Appendix~\ref{app:loop-compaction-metrics} gives the definitions and run coverage.

% RQ3 trace metrics and observed task outcomes.
\begin{table*}[t]
  \centering
  \begin{minipage}[t]{0.58\textwidth}
    \vspace{0pt}
    \centering
    \captionsetup{skip=5pt}
    \caption{Loop trace metrics.}
    \label{tab:rq3-loop-engineering}
    \footnotesize
    \setlength{\tabcolsep}{4pt}
    \renewcommand{\arraystretch}{1.18}
    \begin{tabular*}{\linewidth}{@{\extracolsep{\fill}}l c c c@{}}
      \toprule
      \textbf{Loop Implementation} & \textbf{Rounds} & \textbf{Reg/run} & \textbf{RR (\%)} \\
      \midrule
      Codex goal mode & 32.76 & 0.34 & 20.59 \\
      Claude Code goal mode & 34.69 & 0.13 & 17.65 \\
      Claude dynamic workflows & 97.96 & 0.36 & \textbf{24.11} \\
      Ralph loop & 13.24 & 0.17 & 7.84 \\
      \bottomrule
    \end{tabular*}
  \end{minipage}\hfill
  \begin{minipage}[t]{0.38\textwidth}
    \vspace{0pt}
    \centering
    \captionsetup{skip=5pt}
    \captionof{figure}{Compaction sensitivity.}
    \label{fig:rq3-compaction-ablation}
    \includegraphics[width=\linewidth]{figures/assets/rq3-compaction-ablation.pdf}
  \end{minipage}
\end{table*}


\textbf{(1) Context renewal differs across loop implementations.}
After counting each compaction interval as a round, dynamic workflows record 97.96 context-budget rounds per run, compared with 34.69 for Claude Code goal mode, 32.76 for Codex goal mode, and 13.24 for Ralph. The dynamic workflow run has the highest Resolve Rate at 24.11\%, while Ralph reaches 7.84\%. These values describe the evaluated configurations.
\Needspace{4\baselineskip}
\textbf{(2) Context-budget renewal does not remove regression pressure.}
Dynamic workflows record 0.36 regression events per run despite distributing work across narrower worker contexts. Codex goal mode records 0.34, Claude Code goal mode 0.13, and Ralph 0.17. Dynamic routing therefore reduces dependence on a single growing transcript, but completed obligations still require explicit state tracking when workers return. Appendix~\ref{app:rq3-trace-scope} states the trace scope for worker local contexts.
\begin{summary}
\textbf{RQ3:} Loop engineering changes how context is renewed and how residual work is routed. Dynamic workflows achieve the strongest observed outcome, while regression events remain visible across all four loop profiles.
\end{summary}
